## AUTORES
- (HECHO) Autores sin números.

## KEYWORDS
- (HECHO) Keywords in mayúscula.

## (HECHO) INTRODUCTION

- (HECHO) Me parece que la introducción es demasiado directa. Primero habría que dejar más claro qué es AutoML, HGML, etc. Y después los retos. Quizás se mezclan cosas y la introducción debería quedar:
1º Contextualizar y explicar conceptos

2. Limitaciones o retos que se presentan

3. Cómo se va a abordar esos retos

No mezclar en un párrafo los tres elementos previos.

- (HECHO) Quitar protocolos de Introduction.

## (Cambios hechos pero REVISAR) BACKGROUND

- Me parece que habría que cambiar la estructura de esto. Quizás la nueva estructura pueda ser:

A. AutoML systems (y aquí hablar todo lo que se necesite de Ludwig)

B. HGML/HCAI

Antes de escribir el apartado A, poner una pequeña introducción


## (HECHO, REVISAR) APPLICATION

- Párrafo introductorio. Mover explicación de HAMELIN a otra parte.

- Participatory Design Process
Quizás aquí lo que pondría es cómo se da soporte a tres tipo de estudios. Es decir, adaptar lo siguiente que ahora está en II A: Clinical research employs distinct frameworks or protocols for their research, each with specific design principles and
analytical constraints [5]. Patient registries collect long-term data from people who share the same medical condition. These
registries are typically observational and aim to understand disease history, treatment patterns, and long-term outcomes.
Observational studies test specific hypotheses about exposure-outcome associations without intervention, requiring careful
adjustment for confounding variables and selection bias. Clinical trials evaluate interventions through randomization, blinding
and rigorous protocols defining inclusion/exclusion criteria, temporal phases (screening, baseline, intervention, follow-up,
endpoint evaluation) and endpoint hierarchies (primary endpoints for efficacy, secondary endpoints for safety and exploratory
outcomes).
These distinctions impose specific requirements on predictive modeling that generic ML pipelines ignore [6]. For instance,
in a clinical trial, variables collected during screening (e.g., blood pressure to determine eligibility) have different meanings
than the same variables collected at follow-up (e.g., blood pressure as an outcome). A model trained without awareness of
these temporal and semantic constraints risks producing invalid results, such as data leakage, using information collected after
the outcome to predict that outcome, or confounding, failing to account for variables that influence both exposure and outcome.

- (HECHO) Se podría borrar imagen de mockup.


## (HECHO) ALIGNMENT WITH HUMAN-GUIDED ML

Lo suyo sería que la evaluación sea más global y que se ha verificado que la herramienta proporciona soporte para Four real clinical protocols were served as use cases to ground the interface design in authentic clinical research scenarios:
one patient registry, one observational study, and two clinical trials. These are the details of each protocol that were considered
to ensure the interface was tailored as closely as possible to the required needs.
• EPA-Registro COVID-19: A multicenter patient registry documenting clinical characteristics, treatments, and outcomes
of COVID-19 patients. The registry collected demographic variables (age, sex), comorbidities (hypertension, diabetes,
cardiovascular disease), clinical variables (symptoms, laboratory values), treatments (antivirals, corticosteroids, oxygen
therapy), and outcomes (hospital discharge, ICU admission, death).
• UPPS (Pressure Injuries in Neurosurgery): An observational study investigating the incidence and risk factors of
pressure injuries in neurosurgery patients undergoing procedures longer than 3 hours. The protocol specified variables
related to patient characteristics (age, sex, BMI, Braden scale score), intraoperative factors (surgical duration, patient
positioning, use of pressure-relief surfaces, intraoperative temperature), and postoperative outcomes (presence/absence of
pressure injuries, injury severity).
• COLCHI-COVID: A randomized, double-blind, placebo-controlled clinical trial evaluating the efficacy and safety of
colchicine in hospitalized COVID-19 patients. The protocol defined primary endpoints (clinical improvement at day 14),
secondary endpoints (mortality, ICU admission, adverse events), and detailed inclusion/exclusion criteria (age ≥ 18 years,
confirmed COVID-19, no contraindications to colchicine).
• VACAB-2022: A randomized clinical trial evaluating the efficacy of mucosal bacterial vaccines (autovaccines) compared
to suppressive antibiotic therapy in patients with hip or knee prosthesis infections. The protocol defined primary endpoints
(infection recurrence at 12 months), secondary endpoints (quality of life, adverse events), and inclusion/exclusion criteria
(confirmed prosthesis infection, no immunosuppression, no allergy to vaccine components).